% Generated by Listas/make_statements.py from pratica-lecture-06.md, item 14. Do not edit.
\textbf{Source.} Gibbons (1992), \emph{Game Theory for Applied Economists}, Ch. 2 Problems, Section 2.3, Problem 2.15, p.~136 {[}PDF 143{]}. Reproduced verbatim. Answer in the source: none (Gibbons prints no answers).



\begin{quote}

2.15. Suppose there are \(n\) firms in a Cournot oligopoly. Inverse demand is given by \(P(Q) = a - Q\), where \(Q = q_1 + \cdots + q_n\). Consider the infinitely repeated game based on this stage game. What is the lowest value of \(\delta\) such that the firms can use trigger strategies to sustain the monopoly output level in a subgame-perfect Nash equilibrium? How does the answer vary with \(n\), and why? If \(\delta\) is too small for the firms to use trigger strategies to sustain the monopoly output, what is the most-profitable symmetric subgame-perfect Nash equilibrium that can be sustained using trigger strategies?

\end{quote}

